\section{Metodología de simulación}

\subsection{Motor de eventos discretos}

El núcleo se implementó en Python con la biblioteca estándar, sin SimPy. El
motor mantiene una cola de eventos con \texttt{heapq} y procesa cada evento en
orden de tiempo; los empates se resuelven por prioridad y, finalmente, por
secuencia FIFO. El reloj avanza al instante del siguiente evento, de modo que
no se descompone el tiempo en pasos fijos. El motor está separado de los
modelos: recibe callbacks que programan llegadas, salidas y eventos de control.

\subsection{Generadores aleatorios y reproducibilidad}

Cada fuente de aleatoriedad usa un flujo independiente: llegadas Tipo~1,
llegadas Tipo~2, servicios Tipo~1 y servicios Tipo~2. Las réplicas usan
semillas explícitas consecutivas a partir de una semilla base
($1,2,\ldots,N$). Para comparar escenarios se usan \emph{números aleatorios
comunes} (CRN): las réplicas de todos los escenarios comparten las mismas
semillas, lo que reduce la varianza de las diferencias pareadas. Las
distribuciones se muestrean por transformada inversa, lo que permite
reproducir exactamente una corrida con la misma semilla.

\subsection{Diseño experimental}

El M/M/1 se ejecuta con 50 réplicas de 30\,000 minutos y un warm-up de 1\,000
minutos; el piloto de warm-up compara 0, 500, 1\,000, 2\,000 y 5\,000 minutos
con 10 réplicas y una tolerancia del $0{,}5$\,\%. El SOC se ejecuta con 30
réplicas de 30 días de operación tras 1 día de warm-up (1\,440 minutos), es
decir, una ventana medida de 43\,200 minutos por réplica. Al final del
horizonte se drena el sistema: las alertas admitidas que siguen en cola se
terminan de atender, lo que evita censurar sus tiempos de espera y de
resolución. Los resultados principales usan drenaje; la opción \texttt{--no-drain}
permite evaluar la sensibilidad.

\subsection{Estimación e incertidumbre}

Para cada réplica se calculan las métricas de tiempo y de longitud de cola con
acumuladores ponderados por tiempo y con muestras por alerta. Luego se obtiene
la media entre réplicas, la desviación estándar, el error estándar y un
intervalo de confianza del 95\,\% con la distribución $t$ de Student. También
se reporta el error relativo y, para comparaciones, la diferencia pareada
``alternativa menos referencia'' con su intervalo de confianza. Todos los
resultados, figuras y tablas del informe se generan desde los archivos
versionados en \texttt{results/}; ninguna cifra se escribe a mano.
